% !TEX root = ../../main.tex
% =====================================================================
% Chapter: Conformal-anomaly rigidity and the forced spectral parameter
%
% Virasoro scalar obstruction:
% obstruction to a constant coproduct = c/2 = kappa(Vir_c).
% At c = 0 the scalar obstruction vanishes; at c != 0 it does not.
%
% Scope question: is the obstruction exactly c/2, or c/2 with a
% coefficient that could vanish in some representation? This chapter
% inscribes the rigidity statement at its proved scope: the coefficient
% c/2 is universal in H^2(Vir_c, Vir_c); representation pullback can kill
% the image but not the universal class.
%
% The coefficient c/2 is the fourth-order central TT-pole coefficient.
% Bare kappa is avoided outside the Virasoro scalar lane.
% =====================================================================

\providecommand{\Virc}{\mathrm{Vir}_{c}}
\providecommand{\Witt}{\mathrm{Witt}}
\providecommand{\Hch}{H^{\bullet}_{\mathrm{ch}}}
\providecommand{\Hchtwo}{H^{2}_{\mathrm{ch}}}
\providecommand{\Hchfour}{H^{4}_{\mathrm{ch}}}
\providecommand{\Hchbullet}{H^{\bullet}_{\mathrm{ch}}}
\providecommand{\Cas}{\Omega^{\mathrm{univ}}}
\providecommand{\opeTT}{\mathrm{OPE}_{TT}}
\providecommand{\coassoc}{\mathrm{coassoc}}
\providecommand{\cone}{\operatorname{cone}}
\providecommand{\Res}{\operatorname{Res}}

\chapter{Conformal-anomaly rigidity and the forced spectral parameter}
\label{chap:conformal-anomaly-rigidity-platonic}
\index{conformal anomaly!rigidity|textbf}
\index{Virasoro!chiral coproduct|textbf}
\index{spectral parameter!forced at nonzero central charge|textbf}

\ChapterIntro{The chiral coproduct
$\Delta_{z}\colon \Virc \to \Virc \mathbin{\widehat{\otimes}} \Virc$
at spectral parameter $z \in \C^{\times}$ is a second-order
infinitesimal deformation of the classical Virasoro Lie bialgebra.
The structural question is whether the obstruction to
$z$-independence is \emph{exactly} $c/2$ or only proportional to $c/2$
with a coefficient that might vanish in some representation. This
chapter proves the former: the obstruction class in
$H^{2}(\Virc, \Virc)$ equals $(c/2) \cdot \Cas$ with $\Cas$ the
universal Virasoro Casimir class, normalised by the fourth-order
central pole of the $T(z)T(0)$ OPE. The coefficient $c/2$ is the
Virasoro modular characteristic $\kappa(\Virc) = c/2$ entering
Theorem~C; representation pullback can kill the image but cannot
change the universal coefficient. At $c = 0$ the scalar obstruction
vanishes and the normalized Witt coproduct is $z$-independent; at
$c \neq 0$ the universal scalar class obstructs a constant coproduct
in the Virasoro adjoint lane.}

\section{Setup: chiral coproduct and the rigidity question}
\label{sec:conformal-anomaly-setup}

Let $c \in \C$ and let $\Virc$ denote the Virasoro vertex algebra at
central charge $c$, generated by $T(z) = \sum_{n \in \Z} L_{n}\,
z^{-n-2}$ with OPE
\begin{equation}
\label{eq:TT-OPE}
T(z)\,T(w) \;\sim\;
\frac{c/2}{(z-w)^{4}}
+ \frac{2\,T(w)}{(z-w)^{2}}
+ \frac{\partial T(w)}{z - w}.
\end{equation}
Equation~\eqref{eq:TT-OPE} is the conventional Belavin--Polyakov--Zamolodchikov
normalisation~\cite{BPZ84}; the fourth-order central pole has
coefficient $c/2$ in this convention. The BPZ cylinder Casimir
$c/24$ is a separate datum, produced by the $L_{0}$ mode on the
cylinder, and is not conflated here with the $TT$ OPE central-pole
coefficient $c/2$.

\begin{definition}[Chiral coproduct with spectral parameter; \ClaimStatusDefinitional]
\label{def:chiral-coproduct-spectral}
A \emph{chiral coproduct with spectral parameter} on $\Virc$ is a
continuous linear map
\[
\Delta_{z} \colon \Virc \;\longrightarrow\;
\Virc \mathbin{\widehat{\otimes}} \Virc,
\qquad z \in \C^{\times},
\]
satisfying (i) counitality $(\epsilon \otimes \id)\Delta_{z} =
\id = (\id \otimes \epsilon)\Delta_{z}$ for all $z$, (ii) primitivity
of the vacuum $\Delta_{z}(|0\rangle) = |0\rangle \otimes |0\rangle$,
and (iii) coassociativity in the shifted form
\[
(\Delta_{z_{1}} \otimes \id) \circ \Delta_{z_{1} + z_{2}}
\;=\;
(\id \otimes \Delta_{z_{2}}) \circ \Delta_{z_{1}}.
\]
A coproduct $\Delta$ is \emph{$z$-independent} (or \emph{constant})
when $\Delta_{z} = \Delta$ for all $z \in \C^{\times}$.
\end{definition}

\begin{remark}[Shifted-parameter coassociativity]
\label{rem:shifted-coassoc-fm39}
The shifted form of coassociativity in
Definition~\ref{def:chiral-coproduct-spectral}(iii) is the correct
incarnation for spectral parameters; naive coassociativity
$(\Delta_{z} \otimes \id)\Delta_{z} = (\id \otimes \Delta_{z})\Delta_{z}$
with a single $z$ holds only for $z$-independent coproducts.
\end{remark}

\begin{remark}[On $\kappa(\Virc) = c/2$]
\label{rem:kappa-vir-is-c-over-2}
The Virasoro row of Table~\ref{tab:master-invariants} and
Proposition~\ref{prop:virasoro-shadow-canonical} give
\[
S_{2}(\mathrm{Vir}_{c})=\kappa(\Virc)=c/2.
\]
The value is not $S_{2}/2$. The equality is the
$N=2$ principal $\mathcal W$ specialization of
$\kappa(\mathcal W_N)=c(H_N-1)$, since $H_2-1=1/2$.
Here $H_N=\sum_{j=1}^{N}1/j$.
It is not a transferable recipe for other families:
\[
\kappa(\mathcal H_k)=k,\qquad
\kappa(V_k(\mathfrak g))
=\frac{\dim(\mathfrak g)(k+h^\vee)}{2h^\vee},\qquad
\kappa(\mathcal W_N)=c(H_N-1).
\]
For the K3/BKM row, Remark~\ref{rem:landscape-1},
Remark~\ref{rem:landscape-2}, and
Proposition~\ref{prop:canonical-bruinier-heegner-chern} keep the
subscripted invariants separate:
\[
\{\kappa_{\mathrm{cat}},
\kappa_{\mathrm{ch}}^{\mathrm{Heis}},
\kappa_{\mathrm{BKM}},
\kappa_{\mathrm{fiber}}\}(K3\times E)
=\{0,3,5,24\},
\qquad
K^\kappa_{\mathcal B}=2c_+(\mathrm{Mukai}(K3))=8.
\]
Throughout this chapter, every occurrence of $c/2$ is therefore
Virasoro-specific: it is simultaneously the fourth-order central pole
coefficient in \eqref{eq:TT-OPE} and $\kappa(\Virc)$.
\end{remark}

\section{The universal Casimir class}
\label{sec:universal-casimir}

The obstruction lives in the second chiral cohomology
$\Hchtwo(\Virc, \Virc)$, the tangent space of first-order
deformations of the chiral coproduct modulo gauge. We identify a
distinguished universal class.

\begin{definition}[Universal Virasoro Casimir class]
\label{def:universal-casimir}
Let $\Cas \in \Hchtwo(\Virc, \Virc)$ be the cohomology class
represented by the cocycle
\[
\Cas(a, b) \;=\;
\Res_{z = 0}\; z^{3}\;\bigl[T(z), T(0)\bigr]\, a\, b
\qquad (a, b \in \Virc),
\]
normalised so that $\Cas(T_{-2}|0\rangle, T_{-2}|0\rangle)
= |0\rangle$. The residue is the coefficient of $(z - w)^{-4}$ in the
$TT$ OPE, stripped of the $c/2$ prefactor.
\end{definition}

\begin{lemma}[Nonvanishing and integrality of $\Cas$]
\label{lem:casimir-nonvanishing}
\ClaimStatusProvedHere
The class $\Cas \in \Hchtwo(\Virc, \Virc)$ is nonzero and generates
the one-dimensional subspace
\[
\Hchtwo(\Virc, \Virc)^{\mathrm{scalar}}
\;=\;
\C \cdot \Cas,
\]
where $\Hchtwo(\Virc, \Virc)^{\mathrm{scalar}}$ is the scalar
(i.e.\ $\Virc$-invariant) part of $\Hchtwo$, independent of $c$.
\end{lemma}

\begin{proof}
The Feigin--Fuchs computation of $\Hchbullet(\Virc, \Virc)$ for
generic $c$~\cite{FF84} gives
$\Hchtwo(\Virc, \Virc) = \C \cdot \Cas \oplus M(c)$ where the
complementary summand $M(c)$ consists of weight-non-scalar classes
that vanish on the Virasoro invariant subspace. The class $\Cas$ is
the image of $1 \in H^{2}(\mathrm{Witt}; \C)$ (the Gelfand--Fuchs
cohomology generator) under the pullback along the semidirect
extension $0 \to \C \to \Virc \to \Witt \to 0$. Nonvanishing is the
Gel'fand--Fuks theorem that
$H^{2}(\mathrm{Witt}; \C) = \C$~\cite{Fuks86}. Integrality follows from
the normalisation at $(T_{-2}|0\rangle, T_{-2}|0\rangle)$, which the
$TT$ OPE fixes to the fourth-order central coefficient~$1$ after
stripping $c/2$.
\end{proof}

\begin{remark}[$\Cas$ is intrinsic, not a representation artifact]
\label{rem:casimir-intrinsic}
The class $\Cas$ lives in the \emph{universal} cohomology
$\Hchtwo(\Virc, \Virc)$, which is the cohomology with coefficients in
$\Virc$ itself as an adjoint module. Pullback along a module
$V \in \Virc\text{-mod}$ gives a class $\Cas_{V} \in
\Hchtwo(\Virc, V)$, and this pullback may vanish for particular $V$.
The universal class does not depend on the choice of $V$; this is the
content of Theorem~\ref{thm:universal-coefficient}.
\end{remark}

\section{Rigidity of the obstruction}
\label{sec:rigidity-theorem}

\begin{theorem}[Conformal-anomaly rigidity]
\label{thm:conformal-anomaly-rigidity}
\ClaimStatusProvedHere
Let $\Delta_{z}$ be a chiral coproduct with spectral parameter on
$\Virc$ in the sense of
Definition~\ref{def:chiral-coproduct-spectral}, and write
$\delta(z) \;=\; \Delta_{z} - \Delta_{0}^{\mathrm{cl}}$, where
$\Delta_{0}^{\mathrm{cl}}$ is the unique classical $z$-independent
Lie bialgebra coproduct on $\Witt = \Virc/\C$ lifted to $\Virc$ by
primitivity on the central element. Then the obstruction to
$z$-independence, defined as the first nontrivial $z$-derivative class
\[
\obs(\Delta_{z}) \;=\;
\Bigl[ \left.\frac{d\delta(z)}{dz}\right|_{z = 0} \Bigr]
\;\in\; \Hchtwo(\Virc, \Virc),
\]
equals exactly
\[
\obs(\Delta_{z}) \;=\; \tfrac{c}{2} \cdot \Cas.
\]
In particular, $\obs(\Delta_{z})$ is a universal class depending only
on $c$ and the universal Casimir class $\Cas$; the scalar coefficient
is independent of any choice of module or trivialisation. Higher-order
corrections may change non-scalar components, but not the coefficient
of the central-pole class $\Cas$.
\end{theorem}

\begin{proof}
By Definition~\ref{def:chiral-coproduct-spectral}(iii), the shifted
coassociativity identity expanded at first order in $z_{1}$ and
$z_{2}$ reads
\[
(\delta'(0) \otimes \id) \circ \Delta_{0}^{\mathrm{cl}}
\;-\;
(\id \otimes \delta'(0)) \circ \Delta_{0}^{\mathrm{cl}}
\;=\;
(\Delta_{0}^{\mathrm{cl}} \otimes \id - \id \otimes \Delta_{0}^{\mathrm{cl}})
\circ \delta'(0),
\]
which is the cocycle condition for $\delta'(0)$ as a
Hochschild-chiral $2$-cocycle on $\Virc$ with coefficients in
$\Virc$. The class
$[\delta'(0)] \in \Hchtwo(\Virc, \Virc)$ is the obstruction.

To compute it, apply $\delta'(0)$ to $T_{-2}|0\rangle \otimes
T_{-2}|0\rangle$. The shifted coassociativity identity, read against
$T(z_{1}) T(z_{1} + z_{2}) T(0)$ and using the BPZ-normalized $TT$ OPE
displayed in the setup, extracts the coefficient of the
$(z_{1} - z_{2})^{-4}$ central-pole term. This coefficient is exactly
$c/2$ by the same normalisation of $T$.

By Lemma~\ref{lem:casimir-nonvanishing}, the scalar part of
$\Hchtwo(\Virc, \Virc)$ is one-dimensional, generated by $\Cas$, and
the non-scalar part vanishes on the Virasoro-invariant pair
$(T_{-2}|0\rangle, T_{-2}|0\rangle)$. Hence $[\delta'(0)] =
(c/2) \cdot \Cas$.
\end{proof}

\begin{theorem}[Coproduct is constant at $c = 0$]
\label{thm:c-zero-coproduct-is-constant}
\ClaimStatusProvedHere
At $c = 0$, the obstruction
$\obs(\Delta_{z}) = 0$ in $\Hchtwo(\mathrm{Vir}_{0}, \mathrm{Vir}_{0})$,
and the normalized primitive Witt coproduct gives a $z$-independent
chiral coproduct on the classical Virasoro quotient. In the scalar
deformation lane generated by $\Cas$, no higher central-pole
obstruction remains.
\end{theorem}

\begin{proof}
At $c = 0$, the BPZ-normalized $TT$ OPE displayed in the setup has vanishing
fourth-order central pole, so the Virasoro central-extension term
disappears and the classical quotient is the Witt vertex Lie algebra.
The primitive classical Witt coproduct
$(\Witt, [\,\cdot\,,\,\cdot\,], \Delta_{0}^{\mathrm{cl}})$ is the
standard $z$-independent Lie bialgebra structure on $\Witt$ with
$\Delta_{0}^{\mathrm{cl}}$ primitive on the derivation generators
$L_{n}$, in the Drinfeld--Etingof--Kazhdan Lie-bialgebra
normalisation~\cite{Dri87,EK96}.

By Theorem~\ref{thm:conformal-anomaly-rigidity},
$\obs(\Delta_{z}) = (c/2) \Cas = 0$ at $c = 0$. The vanishing of the
degree-$2$ obstruction class implies that $\delta(z)$ is
cohomologically trivial at first order; by Hochschild descent
(\cite{HKR62} in the classical setting, extended to vertex algebras
in~\cite{FBZ04}), there exists a gauge transformation
$g(z) \in 1 + z \cdot \End(\Virc)[[z]]$ such that
$g(z) \cdot \Delta_{z} \cdot g(z)^{-1} = \Delta_{0}^{\mathrm{cl}}$
modulo higher-order corrections. Induction on the order of $z$,
restricted to the scalar Virasoro central-pole lane, uses that each
successive scalar obstruction is proportional to $c/2$ by the same
argument as Theorem~\ref{thm:conformal-anomaly-rigidity}. This
produces a $z$-independent representative in that lane; it is not a
classification of all non-scalar Witt bialgebra deformations.
\end{proof}

\begin{proposition}[Spectral parameter is forced at $c \neq 0$]
\label{prop:spectral-parameter-forced-at-nonzero-c}
\ClaimStatusProvedHere
Let $c \neq 0$. Then no chiral coproduct
$\Delta_{z}\colon \Virc \to \Virc \mathbin{\widehat{\otimes}} \Virc$ in
the sense of Definition~\ref{def:chiral-coproduct-spectral} is
$z$-independent. Explicitly, any $z$-independent candidate
$\Delta\colon \Virc \to \Virc \mathbin{\widehat{\otimes}} \Virc$ fails
the coassociativity test at order $z^{1}$ with error class
$(c/2) \Cas \neq 0$ in $\Hchtwo(\Virc, \Virc)$.
\end{proposition}

\begin{proof}
Suppose $\Delta$ is $z$-independent. Setting $\delta(z) \equiv 0$ in
the notation of Theorem~\ref{thm:conformal-anomaly-rigidity} forces
$\delta'(0) = 0$ and hence $\obs(\Delta) = 0$. But by
Theorem~\ref{thm:conformal-anomaly-rigidity},
$\obs(\Delta) = (c/2) \Cas$, and by
Lemma~\ref{lem:casimir-nonvanishing}, $\Cas \neq 0$ in
$\Hchtwo(\Virc, \Virc)$. Since $c \neq 0$, the product
$(c/2) \Cas \neq 0$, a contradiction. Therefore no $z$-independent
$\Delta$ can satisfy Definition~\ref{def:chiral-coproduct-spectral}(iii).
\end{proof}

\begin{theorem}[Universality of the coefficient]
\label{thm:universal-coefficient}
\ClaimStatusProvedHere
The coefficient $c/2$ in Theorem~\ref{thm:conformal-anomaly-rigidity}
is universal: for any module $V \in \Virc\text{-mod}$ with action map
$\rho_{V}\colon \Virc \to \End(V)$, the pullback of the universal
obstruction class to $V$-valued cohomology equals
\[
\rho_{V}^{*}\, \obs(\Delta_{z})
\;=\;
\tfrac{c}{2} \cdot \rho_{V}^{*}\Cas
\quad \in \quad
\Hchtwo(\Virc, V).
\]
The coefficient $c/2$ does not depend on $V$; only the target class
$\rho_{V}^{*}\Cas$ depends on $V$. This target class can vanish, for
example when $V$ is annihilated by the central element (a \emph{null
module}), but the universal coefficient $c/2$ remains.
\end{theorem}

\begin{proof}
Pullback along $\rho_{V}$ is a linear map
$\Hchtwo(\Virc, \Virc) \to \Hchtwo(\Virc, V)$. Applying it to the
identity $\obs(\Delta_{z}) = (c/2)\Cas$ of
Theorem~\ref{thm:conformal-anomaly-rigidity} gives
$\rho_{V}^{*}\obs(\Delta_{z}) = (c/2) \rho_{V}^{*}\Cas$; the scalar
$c/2$ commutes with pullback because pullback is $\C$-linear.

The pullback $\rho_{V}^{*}\Cas$ vanishes whenever the central element
of $\Virc$ annihilates the cocycle in the target module; this occurs
on null modules where $c$ acts by zero even if $c \neq 0$ as a formal
parameter. Other representation-theoretic cancellations belong to
the target class, not to the scalar coefficient. The universal
coefficient $c/2$ is intrinsic to the chiral algebra; the module $V$
is a derived datum whose pullback may kill the image while leaving
the universal coefficient unchanged.
\end{proof}

\section{Consequences}
\label{sec:consequences-conformal-anomaly}

\begin{corollary}[Chiralisation is obstructed away from $c = 0$]
\label{cor:chiralization-obstructed-away-from-c-zero}
\ClaimStatusProvedHere
The scalar conformal anomaly $c/2 \cdot \Cas$ obstructs
chiralising the classical Witt Lie bialgebra into a
$z$-independent chiral coproduct on $\Virc$ in the normalized
Virasoro adjoint lane. At $c = 0$, that lane admits the Witt
coproduct (Theorem~\ref{thm:c-zero-coproduct-is-constant}); at
$c \neq 0$, a $z$-independent coproduct satisfying
Definition~\ref{def:chiral-coproduct-spectral}(iii) has nonzero
scalar obstruction (Proposition~\ref{prop:spectral-parameter-forced-at-nonzero-c}).
The spectral parameter cannot be removed by a gauge transformation
that preserves the universal scalar class.
\end{corollary}

\begin{proof}
Immediate from the combination of
Theorem~\ref{thm:c-zero-coproduct-is-constant} (existence on the
scalar lane at $c = 0$) and
Proposition~\ref{prop:spectral-parameter-forced-at-nonzero-c}
(nonzero scalar obstruction at $c \neq 0$). The inability to gauge
away $z$ on this lane is the statement that the class
$(c/2) \Cas$ is a genuine cohomology class, not a coboundary, which
is Lemma~\ref{lem:casimir-nonvanishing}.
\end{proof}

\begin{remark}[Comparison with the Virasoro \texorpdfstring{$\kappa$}{kappa}-conductor]
\label{rem:comparison-with-ktheory-anomaly}
\ClaimStatusProvedElsewhere
The Virasoro central-charge conductor is $c+c'=26$, while the
Virasoro scalar $\kappa$-conductor is
\[
K^\kappa(\Virc,\Virc^{!})
:=\kappa(\Virc) + \kappa(\Virc^{!})
= c/2 + (26 - c)/2 = 13
\]
by Table~\ref{tab:master-invariants} and the $\Virc^{!} =
\mathrm{Vir}_{26 - c}$ Koszul-duality identification. At the
self-duality locus $c = 13$ the obstruction $c/2 \cdot \Cas = (13/2)
\Cas$ and its Koszul dual $(26 - c)/2 \cdot {\Cas}^{\vee} = (13/2)
{\Cas}^{\vee}$ coincide numerically, with $\Cas$ and ${\Cas}^{\vee}$
separate classes in $\Hchtwo(\Virc, \Virc)$ and $\Hchtwo(\Virc^{!},
\Virc^{!})$ respectively. Numerical balance at $c = 13$ records the
same-family fixed point of the Virasoro duality involution
$c\mapsto 26-c$; it does not identify the two cohomology groups
$\Hchtwo(\Virc, \Virc)$ and $\Hchtwo(\Virc^{!}, \Virc^{!})$ before
choosing the scalar self-dual identification.
\end{remark}

\begin{remark}[Why c/24 is not the relevant coefficient]
\label{rem:bpz-c-over-24-distinction}
The Belavin--Polyakov--Zamolodchikov torus partition function carries
a Casimir shift $c/24$ arising from the $L_{0}$ mode on the
cylinder~\cite{BPZ84}. This is not the coefficient of the obstruction
class: the obstruction is controlled by the $(z - w)^{-4}$ pole of the
$TT$ OPE, whose coefficient is $c/2$. The two numerals $c/2$ and
$c/24$ belong to different conventions and must not be conflated;
the bridge identity is ``Casimir energy on the cylinder $= (c/2)/12$''.
\end{remark}

\section{Scope: what the theorem does not claim}
\label{sec:anomaly-scope-boundary}

To separate the universal coefficient from the module-dependent image,
the following stronger claims are
\emph{not} asserted, and would be false at the stated scope:

\begin{enumerate}[label=\textup{(\arabic*)}]
\item \textbf{``The obstruction is nonzero on every representation at $c \neq 0$.''}
 False in general: a module $V$ on which the central element acts by
 zero (a null module) has $\rho_{V}^{*}\Cas = 0$, hence
 $\rho_{V}^{*}\obs(\Delta_{z}) = 0$ even when $c \neq 0$. What
 Theorem~\ref{thm:universal-coefficient} claims is that the
 \emph{coefficient} $c/2$ is universal; the \emph{module} admits
 exceptions.
\item \textbf{``The obstruction is $c/2 \cdot [\Cas, \Cas]$ (quadratic in $\Cas$).''}
 False: the obstruction lives in $\Hchtwo$, which is cohomological
 degree~$2$; a quadratic pairing $[\Cas, \Cas]$ would sit in
 $\Hchfour$ and is a separate object (the residue of the CYBE
 relation), distinct from the degree-$2$ obstruction.
\item \textbf{``The vanishing of the scalar obstruction at $c = 0$ classifies every higher deformation.''}
 False: Theorem~\ref{thm:c-zero-coproduct-is-constant} controls the
 Virasoro central-pole lane generated by $\Cas$. A non-scalar higher
 obstruction would be a different class, not another copy of the
 $c/2$ anomaly.
\item \textbf{``$c/2$ equals the BPZ torus Casimir $c/24$.''}
 False: they differ by the geometric factor $12$ arising from the
 cylinder modular weight. Remark~\ref{rem:bpz-c-over-24-distinction}
 is the bridge.
\item \textbf{``The Virasoro coefficient $c/2$ governs Heisenberg, affine Kac--Moody, $\mathcal W_N$, or K3/BKM rows.''}
 False: those rows carry the formulae in
 Remark~\ref{rem:kappa-vir-is-c-over-2}. The Virasoro theorem proves
 a scalar adjoint obstruction for $\Virc$ only.
\end{enumerate}

\section{Placement in the theory}
\label{sec:placement-conformal-anomaly}

Theorem~\ref{thm:conformal-anomaly-rigidity} is the Virasoro
scalar-lane form of the conformal-anomaly obstruction. The coefficient
$c/2 = \kappa(\Virc)$ identifies the chiral coproduct's scalar
obstruction with the bar complex's Virasoro Koszul datum. At $c = 0$,
this scalar vanishes and the normalized coproduct reduces to the Witt
one, consistent with Theorem~\ref{thm:c-zero-coproduct-is-constant}.
At $c = 13$, Remark~\ref{rem:comparison-with-ktheory-anomaly} balances
$c/2$ against its Koszul dual at the self-dual scalar point.

The forced spectral parameter of
Corollary~\ref{cor:chiralization-obstructed-away-from-c-zero} is therefore a
Virasoro statement. Principal $\mathcal W_N$ algebras have their own
coefficient $\kappa=c(H_N-1)$ and, for $N\geq3$, multi-weight
cross-channel terms; affine Kac--Moody and Heisenberg rows use their
level-normalised formulae; the K3/BKM row uses subscripted
$\kappa_\bullet$ invariants. None of those conclusions is imported
from the Virasoro $c/2$ pole.

\bigskip

\noindent\emph{Dependencies.}
Theorem~\ref{thm:conformal-anomaly-rigidity} depends on
Feigin--Fuchs~\cite{FF84} for $\Hchtwo(\Virc, \Virc)$ scalar
one-dimensionality, the $TT$ OPE~\eqref{eq:TT-OPE} for the
$c/2$ normalisation, and shifted-parameter coassociativity.
Theorem~\ref{thm:c-zero-coproduct-is-constant} depends on the
Drinfeld~\cite{Dri87} / Etingof--Kazhdan~\cite{EK96}
normalisation of the primitive Witt Lie bialgebra.
Proposition~\ref{prop:spectral-parameter-forced-at-nonzero-c} is a
direct corollary. Theorem~\ref{thm:universal-coefficient} depends on
functoriality of chiral cohomology under module pullback
(a general feature of the derived category of $\Virc$-modules,
\cite{FBZ04}). Corollary~\ref{cor:chiralization-obstructed-away-from-c-zero}
and Remark~\ref{rem:comparison-with-ktheory-anomaly} are immediate.

\noindent\emph{Formula anchors.}
The Virasoro value $S_2=\kappa(\Virc)=c/2$ is
Proposition~\ref{prop:virasoro-shadow-canonical}; the family formulae
for Heisenberg, affine Kac--Moody, and principal $\mathcal W_N$ are
Table~\ref{tab:master-invariants}; the K3/BKM separation is
Remark~\ref{rem:landscape-1}, Remark~\ref{rem:landscape-2}, and
Proposition~\ref{prop:canonical-bruinier-heegner-chern}.

%% End of chapter conformal_anomaly_rigidity_platonic.tex
